\section{Conclusion}
\label{sec:conclusion}

This paper presented Tri-Level Hybrid DDQN-ALNS, a learning-augmented optimization framework for the Vehicle Routing Problem with Time Windows (VRPTW). By decomposing metaheuristic search management into a hierarchical decision structure, our method addresses key limitations of both classical stateless heuristics and pure end-to-end neural solvers:
\begin{enumerate}
    \item \textbf{Hierarchical Decision Architecture}: Decoupling macro-level search exploration regimes (Plateau Controller) from micro-level destroy--repair operator selection (Dueling DDQN with Softmax-Entropy confidence gating) and learned acceptance control (LAC) enables coordinated plateau escape while preserving the auditability of neighborhood operators.
    \item \textbf{Strict Feasibility \& Lexicographic Alignment}: Hard temporal and capacity constraints are strictly enforced through $O(1)$ Savelsbergh forward time slack checks and granular spatiotemporal filtering, while multi-objective rewards and exact Set Partitioning recombination over Route Pool $\mathcal{P}$ preserve the primary-secondary fleet minimization hierarchy ($\NV \succ \TD$).
    \item \textbf{Comprehensive Empirical Validation}: Extensive benchmark evaluations across 74 benchmark instances (Solomon-100, Homberger-200, and Homberger-400) under strictly isolated independent cold-starts demonstrate that our method significantly reduces fleet size ($p = 6.91 \times 10^{-7}$) and achieves a $3.73\%$ travel distance reduction under matched fleet counts with zero losses ($p = 9.31 \times 10^{-8}$), passing Bonferroni-corrected significance thresholds.
\end{enumerate}

\subsection{Trade-offs and Limitations}
\label{sec:concl_tradeoffs}
Our empirical findings also delineate the operational scope of the method. The framework is designed as a quality-maximizing optimization pipeline that trades additional computational time for enhanced solution quality, rather than a sub-second search accelerator. On small-scale instances ($N \le 50$), classical heuristics remain computationally preferable. Furthermore, the recombination tier relies on an exact MIP solver, and performance on extreme out-of-distribution topologies depends on policy calibration.

\subsection{Future Research Directions}
\label{sec:concl_future}

Several promising research directions emerge from this work:
\begin{itemize}
    \item \textbf{Electric Vehicle Routing (EV-VRPTW)}: Extending the state space and operator portfolio to incorporate non-linear battery charging profiles, energy recovery, and charging station queuing dynamics.
    \item \textbf{Heterogeneous Multi-Depot Operations}: Generalizing the Macro Plateau Controller to manage inter-depot vehicle transfers and multi-depot fleet capacity constraints.
    \item \textbf{Dynamic Real-Time Re-Optimization}: Adapting the lightweight neural operator selection policy for sub-second reactive re-routing under stochastic customer demand and real-time traffic congestion.
    \item \textbf{Suite-Wide Empirical Ablations}: Scaling component leave-one-out and heuristic-rule sensitivity ablations across the full 74-instance testbed to conduct formal multi-factor ANOVA across diverse spatial-temporal topologies.
\end{itemize}
